\paragraph*{Problems}

\begin{enumerate}[label=\thesection.\arabic*.,ref=\thesection.\arabic*,start=1,series=book-problems]

\item\label{ex:norm-induced-metric} 设 $(X,\|\cdot\|)$ 为赋范向量空间. 证明函数 $d(x,y)=\|x-y\|$ 定义了一个度量.

\item\label{ex:completion-theorem-proof-details} 补全 Theorem~\ref{thm:normed-space-completion} 证明的细节.

\item\label{ex:bounded-continuous-derivatives-banach} 设 $\Omega$ 为 $\mathbb{R}^m$ 中的开集, 用 $C_b^k(\overline\Omega)$ 表示 $\overline\Omega$ 上所有直到 $k$ 阶导数都连续且有界的函数的集合. 定义
\begin{equation}\label{eq:bounded-continuous-derivatives-sum-norm}
\|u\|=\sum_{|\alpha|\leq k}\sup_{\bm{x}\in\overline\Omega}|D^\alpha u(\bm{x})|.
\end{equation}
证明 $C_b^k(\overline\Omega)$ 是 Banach 空间.

\item\label{ex:young-conjugate-exponents-inequality} 证明 Lemma~\ref{lem:young-conjugate-exponents-scalar-inequality}.

\item\label{ex:lp-product-conjugate-exponent-variant} 设 $f\in\widetilde L^p(\Omega)$, $g\in\widetilde L^q(\Omega)$, 其中 $p,q\in(1,\infty)$ 且 $r^{-1}=p^{-1}+q^{-1}<1$. 证明 $fg\in\widetilde L^r(\Omega)$.

\item\label{ex:continuous-lp-not-complete} 证明空间 $\widetilde L^p(\Omega)$ 不完备.

\item\label{ex:bounded-domain-lq-lp-inclusion} 设 $\Omega$ 有界且 $1\leq p<q<\infty$. 证明 $L^q(\Omega)\subset L^p(\Omega)$, 并且存在只依赖于 $\Omega$ 的常数 $C$, 使对所有 $u\in L^q(\Omega)$ 有 $\|u\|_p\leq C\|u\|_q$.

\item\label{ex:lp-two-product-epsilon-bound} 证明对每个 $u,v\in L^2(\Omega)$ 和每个 $\epsilon>0$, 有 \[\left|\int_\Omega uv\,\mathrm{d}\bm{x}\right|\leq\epsilon\|u\|_2^2+C(\epsilon)\|v\|_2^2,\] 其中 $C(\epsilon):=(4\epsilon)^{-1}$.

\item\label{ex:test-functions-dense-lp} 证明 $\mathcal{D}(\Omega)$ 在 $L^p(\Omega)$ 中稠密, $1\leq p<\infty$.

\item\label{ex:inner-product-joint-continuity} 设 $H$ 为内积空间. 证明内积在 $H\times H$ 上连续.

\item\label{ex:sequence-lp-two-cauchy-schwarz-triangle} 写出 $\ell^2$ 和 $L^2(\Omega)$ 中 Cauchy-Schwarz 不等式及三角不等式的具体形式.

\item\label{ex:lp-two-separable} 证明 $L^2(\Omega)$ 可分. Hint: 使用 Problem~\ref{ex:test-functions-dense-lp}.

\item\label{ex:closed-subspace-double-orthogonal-complement} 证明若 $M$ 闭, 则 $(M^\perp)^\perp=M$.

\item\label{ex:finite-dimensional-norms-equivalence} 证明 $\mathbb{R}^n$ 上的所有范数都等价.

\end{enumerate}
